% Pursuing Stacks (cotes 134-2 to 134-8): archive cote 134-7, site page 26 (PDF page 27), typescript page 468.
% typescript. Draft compared: ps6.tex:1208-1266.
% AI-transcribed by Claude Opus 5.5 (claude-opus-5-5, adaptive thinking, effort high), Message Batches API, batch ps3 (msgbatch_01XdFmuVVBAcgc7EFxdtCqsj), 2026-11-12; not yet checked by eye.
\scanpage{134-7/26}
\mspage{468}

abelianization
\[ \text{(C')} \qquad X_{*} \longrightarrow W(L(X_{*})) \quad . \]
Postnikovs construction, for \deleted{\unsure{n-}}\inserted{$(n-1)$-}connected $X_{*}$ , consists in composing
this map with the "augmentation"
\[ W(L(X_{*})) \longrightarrow K(\hand{$\pi$}_n,\hand{$n$}) \deleted{\unsure{.}} \]
\hand{(}where $\hand{$\pi$}_n\hand{$\simeq \mathrm{H}_n$}$ is the first possibly non trivial homology module of the
chain complex corresponding to $L(X_{*})$\hand{)}, and after this \inserted{only} take the homotopy
fiber, and iterate (the homotopy fiber will be n-connected now). This
process, in a way, breaks the natural abelianizations \xout{\illegible{1}} into
pieces, a \xout{\illegible{1}} brutal thing to do one will admit, all the more so as
we start with a beautiful complex with projective components, and kind
of destroy it's unmarred harmony by tearing out or breaking off the most
showy part, $\mathrm{H}_n$ to name it, which now looks so lost and awkward we reall\hand{y}
can't just leave it as it is, we first have to take a projective resolu-%
tion of it, choosing it as we may ... \inserted{But} \unsure{We|we} wo'nt do \inserted{all} this, will we, \over{\unsure{and}}{and}
rather keep abelianization and \hand{(}C'\hand{)} as God gave them to us, and take the
homotopy fiber (\unsure{We|we} hope God will give this too ...), and repeat the
process, without even having to care at any stage which $\hand{$\pi$}_i$\hand{'s} vanish and
which not. Let's call this the "\uline{soft Postnikov dévissage}" , in contrast
to the"brutal"one. In describing the process (which of course makes a
sense in the discrete context as well as in the schematic one), I impli-%
citly admitted that the $X_{*}$ we start with is 1-connected, or for the
very least has abelian $\hand{$\pi$}_1$ (a notion we'll have to come back to, in the
schematic \xout{\illegible{1}} set up). But we may as well apply it to a (discrete)
$K(G,1)$ type, $G$ any discrete group, then it amounts to tak\over{e}{ing} the descen-%
ding filtration of $G$ by iterated commutator groups, which is a finite
filtration iff $G$ is solvable. Maybe it would be more natural still to
take the similar descending filtration, suitable for the study of \uline{nilpo-}%
\uline{tent} groups rather than solvable one's, with
\margin{(D)}
\[ G^{(n+1)} = \hand{$[$}G,\ G^{(\unsure{n+1|n-1})}\hand{$]$}, \]
where \hand{$[$}A,B\hand{$]$} denotes the subgroup of $G$ generated by commutators \xout{\illegible{1}}
\[ (a,b) = aba^{-1}b^{-1}\hand{,} \]
with a in A, b in B. It does'nt seem there is a similar distinction to
make in case we start with a \xout{simply} 1-connected $X_{*}$ , more specifically
if $X_0=X_1= e$. There may be some \inserted{extra} caution needed however when we assume
only $X_0= e$ without assuming 1-connectedness, even when $\hand{$\pi$}_1$ is abelian,

\margin{26}
